\chapter{Reproducibilidad y estructura del repositorio}

\section{Estructura}
\begin{lstlisting}[language=bash]
PINN-HRF-fMRI/
|-- configs/          # hiperparametros y limites
|-- data/             # documentacion y derivados livianos
|-- docs/             # evolucion metodologica
|-- notebooks/        # trece etapas reproducibles
|-- report/           # LaTeX, figuras y PDF
|-- results/          # metricas y figuras finales
|-- scripts/          # auditoria y utilidades
|-- src/pinn_hrf/     # modelo ODE, metricas y estadistica
|-- tests/            # verificaciones fisiologicas basicas
|-- requirements.txt
|-- pyproject.toml
`-- README.md
\end{lstlisting}

\section{Semillas y precisión}
La generación sintética utilizó una semilla base documentada y derivaciones deterministas por escenario, SNR y réplica. Las MLP y PINN almacenaron la semilla de cada experimento. PyTorch operó en \texttt{float64} para reducir errores en derivadas y optimización de parámetros.

\section{Resultados guardados}
Los lotes exportaron métricas por corrida, parámetros, HRF, predicciones, estados de ejecución y resúmenes agregados. La etapa final generó tablas CSV, LaTeX y figuras para evitar transcribir manualmente valores al informe.

\section{Pruebas}
\begin{lstlisting}[language=bash]
python -m pytest
# Resultado verificado: 2 passed
\end{lstlisting}
Las pruebas actuales comprueban homeostasis sin estímulo y respuesta BOLD positiva. Trabajo futuro debe ampliar cobertura a dimensiones, acotación de parámetros, comparación con soluciones analíticas simples y estabilidad numérica.

\section{Auditoría antes de publicar}
Un script recorre el repositorio, clasifica archivos, detecta binarios y marca elementos mayores a 20 MB. La auditoría final no encontró datos crudos ni archivos de gran tamaño destinados a GitHub.
